%!TEX encoding = UTF8
%!TEX root = 0-notes.tex

\unchapter{Limites}
\fancyhead[C]{\textbf{Chapitre \thechapter~— Limites}}
\addcontentsline{toc}{section}{Méthodologie}

placeholder

% interprétation numérique

% asymptotes et interprétation graphique

% limites simples f(x) -> f(a)

% limites de polynômes frac p(x) / q(x)

% fonctions cos/sin

% fonction arctan

% fonction exp/ln/log
% domination des polynômes

\newpage
\addcontentsline{toc}{section}{Exercices}
\fancyhead[C]{\textbf{Chapitre \thechapter~— Exercices}}




\exe{,difficulty=1}{
	Le but de cet exercice est d'étudier le comportement de l'écart-type d'une série statistique lorsqu'une valeur moyenne est ajoutée à répétition.
	\begin{enumerate}
		\item
		Montrer qu'ajouter une valeur moyenne à une série statistique ne change pas sa moyenne.
		\item
		\begin{multicols}{2}
		%En déduire, sans calcul littéral, la moyenne de la série statistique ci-contre, où $N\in\N$ est un entier naturel.
		Donner la moyenne de la série statistique ci-contre, où $N\in\N$ est un entier naturel.
		
		\begin{center}
		\begin{tabular}{|c|c|c|c|}\hline
		Valeur   & 0 & 20 & 10 \\ \hline
		Effectif & 17 & 17 & $N$ \\ \hline
		\end{tabular}
		\end{center}
		\end{multicols}
	
		\item
		Exprimer l'écart-type $\sigma$ de cette série comme fonction de $N$.
		Calculer $\sigma$ lorsque $N$ vaut 10 ; 20 ; 50 ; 100 à l'aide de la fonction.
		%La série devient-elle plus ou moins concentrée autour de sa moyenne ?
		Vers quelle valeur l'écart-type $\sigma$ s'approche-t-il \mbox{lorsque $N$ grandit ?}
		
		\item
		Donner le premier entier $N\in\N$ à partir duquel l'écart-type est inférieur ou égal à 1.
		\item
		Est-il possible, pour n'importe quel seuil $S>0$, de trouver un entier $N\in\N$ tel que l'écart-type soit inférieur à $S$ ?
	\end{enumerate}
}{exe:variance-exemple}{
	%À rendre pour possibilité de points bonus...
}




\newpage
\addcontentsline{toc}{section}{Solutions}
\fancyhead[C]{\textbf{Chapitre \thechapter~— Solutions}}

\shipoutAnswer



